\chapter{Spatial and Ambient Interfaces}

\section{Breaking the Tyranny of the Screen}
Legacy interfaces trap human attention within glowing rectangles, divorcing the citizen from their physical environment. The Sovereign Stack operates in physical reality (Layer 1 physics, Layer 2 hardware). Therefore, the Layer 6 interface must project outward into the spatial world.

\section{Ambient E-Ink and Shared Context}
Not all information requires a personal smartphone. For communal resources (such as the Commons Kitchen in Scenario Gamma or the Apprenticeship Foundry in Scenario Kappa), Layer 6 leverages ambient, low-power e-ink displays mounted directly onto the physical architecture. 

These displays render the current cryptographic state of the room (e.g., current acoustic budget, thermal limits, occupancy rights) into the physical space. Because e-ink only consumes power during a refresh cycle, it provides a persistent, zero-energy UI that naturally integrates into the community's periphery without demanding constant attention.

\section{Spatial AR and the Voxel Grid}
For high-density coordination (like Algorithmic Architecture in Scenario Xi), citizens utilize Augmented Reality (AR) overlays. The Oasis Engine maps cryptographic policy (Layer 5) directly onto the 3D voxel grid of the physical world. A citizen looking at a vacant lot through AR hardware sees the exact zoning permissions, thermal boundaries, and Trust Ring attestations floating as semantic metadata anchored to the physical coordinates. This Spatial UI ensures that the cognitive map of the community perfectly aligns with the physical reality.

\section{Iterative Refinement: Multi-User AR Collision Resolution}
Scenario Xi (Architecture) highlighted an immediate issue: when multiple users occupy the same physical space wearing AR glasses, overlapping digital intents create visual chaos. Layer 6 resolves this via \textit{Multi-User Spatial Collision Resolution}. The local mesh dynamically prioritizes semantic rendering based on the user's current task vector; if two users are looking at the same structural joint, the Engine mathematically degrades conflicting UI elements into subtle boundary lines, ensuring shared spatial coherence without visual overwhelm.

\section{Iterative Refinement: Spatial Audio Zoning}
Communal spaces (Scenarios Gamma and Epsilon) proved that even ambient e-ink can be too visually demanding. The interface layer now explicitly supports \textit{Ambient Non-Visual Cues} and \textit{Spatial Audio Zoning}. By utilizing directional audio parameters, the Oasis Engine can project subtle acoustic boundaries or olfactory triggers (via HVAC integration) to indicate state changes—such as an impending thermal limit—allowing citizens to perceive cybernetic alerts entirely subconsciously, bypassing the visual cortex altogether.
